\documentclass[11pt,aspectratio=169]{beamer}
\usetheme{Madrid}
\usecolortheme{default}
\usepackage[utf8]{inputenc}
\usepackage{booktabs}
\usepackage{graphicx}
\usepackage{array}
\usepackage{tikz}
\usepackage{tcolorbox}
\usepackage{colortbl}
\usetikzlibrary{shadows.blur,positioning}
\definecolor{medixPrimary}{RGB}{0,60,120}
\definecolor{medixAccent}{RGB}{0,185,170}
\definecolor{medixDark}{RGB}{15,20,30}
\definecolor{medixLight}{RGB}{245,248,252}
\setbeamercolor{background canvas}{bg=medixLight}
\setbeamercolor{title}{fg=white}
\setbeamercolor{frametitle}{fg=white}
\setbeamercolor{author in head/foot}{bg=medixDark, fg=white}
\setbeamercolor{date in head/foot}{bg=medixAccent, fg=white}
\setbeamertemplate{navigation symbols}{}
\setbeamersize{text margin left=1.1cm,text margin right=0.7cm}
\setbeamertemplate{frametitle}
{
    \vspace*{-0.12cm}
    \begin{beamercolorbox}[
        wd=\paperwidth,
        ht=1.05cm,
        dp=0ex
    ]{frametitle}
        \begin{tikzpicture}[remember picture,overlay]
            \fill[medixPrimary]
                (current page.north west)
                rectangle
                ([yshift=-1.05cm]current page.north east);
            \fill[medixAccent]
                ([xshift=-2.4cm]current page.north east)
                --
                (current page.north east)
                --
                ([yshift=-1.05cm]current page.north east)
                -- cycle;
            \draw[
                white,
                opacity=0.28,
                line width=0.8pt
            ]
                ([yshift=-0.86cm]current page.north west)
                --
                ([yshift=-0.86cm]current page.north east);
            \node[
                anchor=west,
                white,
                font=\bfseries\Large
            ]
            at ([xshift=1cm,yshift=-0.52cm]current page.north west)
            {\insertframetitle};
        \end{tikzpicture}
    \end{beamercolorbox}
    \vspace{0.42cm}
}
\setbeamertemplate{footline}
{
    \leavevmode
    \hbox{
        \begin{beamercolorbox}[
            wd=.80\paperwidth,
            ht=3.2ex,
            dp=1.4ex,
            left
        ]{author in head/foot}
            \hspace{0.5cm}
            \small\textbf{MedixAI}
        \end{beamercolorbox}
        \begin{beamercolorbox}[
            wd=.20\paperwidth,
            ht=3.2ex,
            dp=1.4ex,
            right
        ]{date in head/foot}
            \small
            \insertframenumber/\inserttotalframenumber
            \hspace{0.45cm}
        \end{beamercolorbox}
    }
}
\setbeamertemplate{footline}
{
    \leavevmode
    \hbox{
        \begin{beamercolorbox}[
            wd=.80\paperwidth,
            ht=3.2ex,
            dp=1.4ex,
            left
        ]{author in head/foot}
            \hspace{0.5cm}
            \small\textbf{MedixAI}
        \end{beamercolorbox}
        \begin{beamercolorbox}[
            wd=.20\paperwidth,
            ht=3.2ex,
            dp=1.4ex,
            right
        ]{date in head/foot}
            \small
            \insertframenumber/\inserttotalframenumber
            \hspace{0.45cm}
        \end{beamercolorbox}

    }
}
\addtobeamertemplate{background}{}
{
    \begin{tikzpicture}[remember picture,overlay]
    \fill[medixPrimary!5]
        (current page.north west)
        rectangle
        ([xshift=0.55cm]current page.south west);
    \draw[
        medixAccent,
        line width=1.6pt
    ]
    ([xshift=0.55cm]current page.north west)
    --
    ([xshift=0.55cm]current page.south west);
    \end{tikzpicture}
}

\title{MediXAI}
\subtitle{
Multimodal AI for Prescription Analysis and Personalized Health Monitoring
}
\author{
23201009 | Momin Ali Rifat \\[0.2cm]
23201023 | Nazmuz Saif \\[0.2cm]
23201033 | Tahsin Rahman Khan
}
\institute{
Department of CSE \\
University of Asia Pacific
}


\date{\today}
\begin{document}
\begin{frame}[plain]
\begin{tikzpicture}[remember picture,overlay]
\fill[medixPrimary]
(current page.south west)
rectangle
(current page.north east);
\fill[white,opacity=0.04]
([xshift=-3cm,yshift=1cm]current page.center)
circle (4cm);
\fill[medixAccent,opacity=0.08]
([xshift=4cm,yshift=-2cm]current page.center)
circle (3cm);
\draw[
medixAccent,
line width=4pt
]
([xshift=1cm,yshift=-1.2cm]current page.north west)
--
([xshift=-1cm,yshift=-1.2cm]current page.north east);
\end{tikzpicture}


\vspace{1.5cm}
\centering
{\color{white}\fontsize{38}{40}\selectfont\bfseries
MediXAI
}

\vspace{0.4cm}
{\color{medixAccent}\fontsize{22}{24}\selectfont
Multimodal AI for Prescription Analysis and Personalized Health Monitoring
}


\vspace{0.15cm}
{\color{white}\large
Department of Computer Science and Engineering
}

\vspace{0.08cm}
{\color{white}\large
University of Asia Pacific
}

\vspace{1cm}
{\color{medixAccent}\large\bfseries Presented By}

\vspace{0.20cm}
{\color{white}\fontsize{12}{14}\selectfont
\begin{tabular}{@{}l c l@{}}
23201009 & | & Momin Ali Rifat \\
23201023 & | & Nazmuz Saif \\
23201050 & | & Tahsin Rahman Khan \\
\end{tabular}
}
\end{frame}




\begin{frame}{Presentation Outline}
\tableofcontents
\end{frame}


\section{Introduction}

\begin{frame}{Introduction}

\textbf{MediXAI Overview}

\vspace{0.3cm}

MediXAI is a multimodal healthcare framework for:
\begin{itemize}
    \item Automated prescription image analysis (handwritten and image)
    \item Personalized health monitoring using patient signals and records
    \item Supporting clinical decision-making systems
\end{itemize}

\vspace{0.4cm}

\textbf{Why it matters?}
\begin{itemize}
    \item Prescription misinterpretation leads to medical risks
    \item Existing systems lack multimodal integration
    \item Demand for AI-based personalized healthcare is increasing
\end{itemize}

\vspace{0.4cm}

\textbf{Key Idea:}
\begin{itemize}
    \item Fusion of image, text, and patient data for intelligent healthcare prediction
\end{itemize}

\end{frame}




\section{Problem Statement}
\begin{frame}{Problem Statement}

\textbf{Major Problem}
\begin{itemize}
    \item Existing systems are single-modal (OCR / text / image only)
    \item No unified multimodal healthcare system exists
\end{itemize}

\vspace{0.3cm}

\textbf{Key Limitations}
\begin{itemize}
    \item High error in real-world prescriptions
    \item Poor performance on noisy / handwritten data
    \item No real-time clinical deployment system
\end{itemize}

\vspace{0.3cm}

\textbf{Motivation}
\begin{itemize}
    \item Prescription misreading causes medication errors
    \item Patients \& pharmacists often misunderstand handwritten notes
    \item Need AI-based safe prescription interpretation system
\end{itemize}

\end{frame}


\section{Objectives of the Review}
\begin{frame}{Objectives of the Review}

\begin{itemize}
    \item \textbf{Model Comparison:} ML/DL, CNN, XAI, OCR, NLP Vision-Language models
    
    \item \textbf{Metrics:} Accuracy, Precision, Recall, F1-score, latency
    
    \item \textbf{Robustness:} Blur, noise, low-quality scans, handwriting variation, multilingual data
    
    \item \textbf{Limitations:} Data scarcity, weak validation, poor generalization, clean-data dependency
    
    \item \textbf{Explainability:} LIME, SHAP, Grad-CAM
\end{itemize}

\vspace{0.4cm}

\textbf{Research Gap:} Separate OCR-based prescription analysis and health monitoring systems; no unified end-to-end clinical AI framework.
\end{frame}





\section{Paper Selection Methodology}

\begin{frame}{Paper Selection Methodology}

\begin{itemize}
    \item \textbf{Total Papers Reviewed:} 21 research papers
    
    \item \textbf{Sources:}
    \begin{itemize}
        \item IEEE Xplore
        \item Google Scholar
        \item Kaggle
        \item ScienceDirect
        \item ResearchGate
    \end{itemize}
    
    \item \textbf{Publication Years:} 2020 -- 2026
    
    \item \textbf{Inclusion Criteria:}
    \begin{itemize}
        \item Medical OCR, prescription analysis, or healthcare AI related studies
        \item ML, DL, multimodal AI, or XAI-based approaches
    \end{itemize}
    \item \textbf{Exclusion Criteria:}
    \begin{itemize}
        \item Non-medical domain papers
        \item Non-AI based traditional rule-based systems
    \end{itemize}
\end{itemize}
\end{frame}





\section{Literature Review}

\begin{frame}{Literature Review}
\small

\centering
\renewcommand{\arraystretch}{1.2}

\begin{tabular}{|m{3.2cm}|m{5.2cm}|m{5cm}|}
\hline

\rowcolor{medixPrimary}
\color{white}\textbf{Category} &
\color{white}\textbf{Methods / Papers} &
\color{white}\textbf{Key Findings}
\\
\hline

OCR + NLP Pipeline  &
OCR + NLP (ViT, BERT), error propagation analysis &
Good extraction (71--84\% accuracy) \\
\hline

CNN + CRNN OCR Systems &
CNN + BiLSTM + CTC loss &
~70--82\% accuracy, affected by imbalance and handwriting variation \\
\hline

Hybrid OCR + NLP  &
CNN-LSTM + Tesseract OCR + NLP &
Around 92\% accuracy, but fails on low-quality images \\
\hline

Hybrid OCR Pipelines &
ANN-based OCR, fax prescription dataset &
CRR 94.5\%, but limited generalization \\
\hline

RAG + OCR Systems &
OCR + NLP + RAG models &
92\% accuracy, but needs human verification \\
\hline

YOLO + OCR &
YOLO + OCR + Spell Correction &
High ROI accuracy (99\%)\\
\hline

\end{tabular}

\vspace{0.2cm}

\textbf{Key Insight:}
Most systems show good accuracy, but real-world clinical deployment is still limited due to data and validation issues.

\end{frame}





\section{Comparison Table}

\begin{frame}{Comparative Analysis}
\vspace{0.3cm}

\scriptsize
\centering

\renewcommand{\arraystretch}{1.2}

\begin{tabular}{|m{0.5cm}|m{3.2cm}|m{1.6cm}|m{3.2cm}|m{3.2cm}|}

\hline

\rowcolor{medixPrimary}
\color{white}\textbf{P} &
\color{white}\textbf{Method} &
\color{white}\textbf{Accuracy} &
\color{white}\textbf{Advantages} &
\color{white}\textbf{Limitations}
\\

\hline

P1 & OCR + NLP Pipeline & 71--84\% & Good system-level analysis of OCR errors & No real clinical dataset \\
\hline

P2 & CNN + CRNN OCR & 70--82\% & Works for basic handwritten recognition & Sensitive to imbalance and handwriting style \\
\hline

P3 & CNN-LSTM + Tesseract OCR & 92\% & Good medical entity extraction & Weak on poor-quality images \\
\hline

P4 & MIRAGE (LLaVA + Qwen-VL + CLIP) & 82\% extraction & Strong multimodal approach & Vision encoder bottleneck \\
\hline

P5 & OCR + NLP + RAG System & 92\% & Structured medical outputs & LLM hallucination issues \\
\hline

P6 & YOLO + OCR + Spell Correction & 96--99\% ROI & Very high detection accuracy & No full end-to-end evaluation \\
\hline

P7 & Donut (OCR-free Transformer) & 84--88\% & End-to-end document understanding & Limited dataset and validation \\
\hline

\end{tabular}

\vspace{0.2cm}

\textbf{Key Insight:} Most methods perform well on specific tasks, but real-world medical prescription understanding is still unreliable.

\end{frame}





\section{Research Trends}
\begin{frame}{Research Trends and Observations}

\footnotesize

\textbf{Popular Methods}
\begin{itemize}
    \item CNN, CNN-LSTM, DONUT, YOLO
    \item OCR + NLP, RAG, Multimodal AI
\end{itemize}

\vspace{0.2cm}

\textbf{Performance Trend}
\begin{itemize}
    \item CNN/RNN: 80--82\%
    \item CNN-LSTM: $\sim$92\%
    \item YOLO: 99.6\% (ROI)
\end{itemize}

\vspace{0.2cm}

\textbf{Key Observations}
\begin{itemize}
    \item Hybrid systems give better real-world performance
    \item Multimodal models are improving understanding
    \item XAI is used for trust and transparency
\end{itemize}

\vspace{0.2cm}

\textbf{Final Insight:}
YOLO performs best for detection, while CNN-LSTM and hybrid OCR+NLP models are most practical for end-to-end medical document analysis.

\end{frame}



\section{Research Gaps}
\begin{frame}{Research Gaps}
\begin{itemize}

\item \textbf{Data Limitation}  
Small and synthetic datasets reduce generalization.

\item \textbf{Deployment Gap}  
Models are not tested in real clinical environments.

\item \textbf{Noise Issue}  
Performance drops on blurred and handwritten prescriptions.

\item \textbf{Integration Gap}  
No unified OCR, NLP, AI, validation system.

\item \textbf{Clinical Validation}  
Limited real hospital or pharmacy testing.

\item \textbf{Interpretability}  
XAI is still not fully reliable for medical decisions.

\end{itemize}

\vspace{0.15cm}

\textbf{Final Insight:}  
Need for robust, real-time, clinically validated healthcare AI systems.

\end{frame}






\section{Future Directions}

\begin{frame}{Future Research Directions}

\begin{itemize}

\item \textbf{Edge AI \& Real-Time Systems}  
Lightweight models for fast on-device prescription processing.

\item \textbf{Federated Learning}  
Privacy-preserving training without central medical data storage.

\item \textbf{Improved Vision Encoders}  
Better handwriting understanding using SigLIP / InternVL.

\item \textbf{Advanced OCR Models}  
Adoption of TrOCR, Donut for robust document extraction.

\item \textbf{Multimodal Healthcare AI}  
Integration of OCR, NLP, and clinical decision support systems.

\item \textbf{Explainable AI Enhancement}  
More reliable XAI using hybrid approaches (LIM, SHAP, CAM).

\end{itemize}
\end{frame}




\section{Conclusion}
\begin{frame}{Conclusion}
\begin{itemize}

\item \textbf{Overall Observation}  
Multimodal OCR systems are effective for medical document understanding but still face real-world challenges.

\item \textbf{Best Approaches}  
Vision-language models with improved encoders (SigLIP, InternVL) and OCR models like DONUT and TrOCR perform better than traditional methods.

\item \textbf{Strengths}  
Good performance on structured data and strong automation capability.

\item \textbf{Limitations}  
Weak performance on noisy handwriting, limited datasets, and lack of clinical validation.

\item \textbf{Final Remark}  
Further improvements are required for reliable real-world healthcare deployment.

\end{itemize}

\end{frame}







\section{Team Contribution}

\begin{frame}{Team Contribution}
\vspace{0.3cm}

\small

\begin{table}

\centering

\renewcommand{\arraystretch}{1.2}

\begin{tabular}{|m{1.5cm}|m{1.5cm}|m{3.5cm}|m{3.5cm}|}

\hline

\rowcolor{medixPrimary}

\color{white}\textbf{Team Member ID} &
\color{white}\textbf{Papers Reviewed} &
\color{white}\textbf{Report Contribution} &
\color{white}\textbf{Presentation Contribution}

\\

\hline

23201009 &
7 Papers &
Introduction, Problem Statement, Objectives, Methodology Literature Review &
Prepared Intro, Methodology and Literature Review Slides

\\

\hline

23201023 &
7 Papers &
Comparison, Research Gaps, Trends &
Prepared Analysis and Comparison Slides

\\

\hline

23201050 &
7 Papers &
Future Work, Conclusion, Discussion, References &
Prepared Final Slides and Editing

\\

\hline

\end{tabular}

\end{table}

\end{frame}








\section{References}

\begin{frame}{References}

\begin{thebibliography}{10}

\bibitem{paper1}
Mankash et al., "MIRAGE: Multimodal Identification and Recognition of Annotations in Indian General Prescriptions," arXiv, 2024.

\bibitem{paper2}
Daniela Gifua, ``AI-backed OCR in Healthcare,'' None, 2023.

\bibitem{paper3}
Shaik Sharjeel and Mohammad Ubaidulla Arif, ``AI-Based OCR System for Handwritten Medical Prescription Recognition and Interpretation,'' None, 2025.

\bibitem{paper4}
Hung Truong Thanh Nguyen et al., `` Evaluation of XAI: SHAP, LIME, and CAM '' Conference Paper, 2021.



\end{thebibliography}

\end{frame}




\begin{frame}[plain]

\begin{tikzpicture}[remember picture,overlay]

\fill[medixPrimary]
(current page.south west)
rectangle
(current page.north east);

\fill[medixAccent,opacity=0.08]
(current page.center)
circle (3cm);

\end{tikzpicture}

\vspace{3cm}

\centering

{\color{white}\fontsize{34}{36}\selectfont\bfseries
Thank You
}

\vspace{0.4cm}

{\color{medixAccent}\Large
Questions?
}

\end{frame}
\end{document}